Ransomware poses an increasing threat to Linux-based servers and virtualized infrastructure, where a single successful encryption process can lead to operational disruption and data loss. Existing Linux ransomware detection approaches often rely on fixed-window feature aggregation or raw continuous features that may generalize poorly to unseen families and evasive partial-encryption behavior.

This thematic paper investigates the feasibility of detecting Linux ransomware through real-time file I/\allowbreak{}O monitoring at the Virtual File System (VFS) layer using the extended Berkeley Packet Filter (eBPF). The prototype combines an eBPF-based VFS monitoring sensor, a compact representation that encodes each file I/\allowbreak{}O event as two 10-bit tokens, and a stateful micro-LSTM (\textmu{}LSTM) classifier for streaming sequence detection.

A quality-audited dataset of 682 labeled VFS-trace recordings was constructed across fourteen ransomware families under ransomware-only, mixed-workload, and benign-only conditions, comprising approximately 598.6 million malware-labeled I/\allowbreak{}O events. Detection was evaluated under a leave-one-family-out (LOFO) protocol. Because nine of the fourteen families were represented by a single binary, their folds reduce to leave-one-binary-out evaluation; the results are therefore interpreted as held-out-family performance within the observed behavioral space rather than universal generalization.

Across the fourteen folds the mean AUC was 0.9752, with ten folds achieving a 0.0\% missed-detection rate (MDR) and eleven below 5\%. Three failure modes were identified --- SougoLock, HelloKitty, and Akira --- each traced to an identified cause. An in-sample heuristic layer reduced false positives from 193 to 4, a result that does not estimate performance on unseen benign workloads. These findings indicate that VFS-layer file I/\allowbreak{}O carries sufficient signal to detect most held-out families, while the identified failures delineate the boundary of what this monitoring layer can observe.

This study contributes: (1) a VFS-level eBPF monitoring methodology and a multi-family labeled Linux ransomware dataset; (2) evidence that the VFS representation itself carries the detection signal; (3) a characterization of what VFS-layer monitoring can and cannot observe; and (4) a compact streaming detector evaluated under cross-family generalization.

CONTENTS

ACKNOWLEDGMENTS	iv

ABSTRACT	v

CHAPTER I	1

INTRODUCTION	1

1.1 Background and Motivation	1

1.2  Problem Statement	3

1.3  Research Objectives	4

1.4  Research Questions	5

1.5  Scope and Research Contributions	6

1.6  Thesis Organization	8

CHAPTER II	10

BACKGROUND AND LITERATURE REVIEW	10

2.1  Ransomware: Characteristics and Behavioral Taxonomy	10

2.1.1  Definition and Threat Landscape	10

2.1.2  Conceptual Ransomware Execution Lifecycle	11

2.1.3  Encryption Strategies and Evasion	12

2.2  eBPF for Kernel-Level Security Monitoring	13

2.2.1  Overview of eBPF	13

2.2.2  eBPF Hook Points and Event Sources	13

2.2.3 VFS-Level Monitoring for File I/\allowbreak{}O Behavior	14

2.2.4 Strengths and Constraints of eBPF Monitoring	14

2.3  Machine Learning for Ransomware Detection	14

2.3.1  Feature-Based and Storage-Layer Approaches	14

2.3.2  Kernel-Level and eBPF-Based Detection	15

2.3.3  Generalization and Deployment Challenges	16

2.4  Datasets for Ransomware Research	19

2.5  Research Gaps and Thesis Positioning	22

CHAPTER III	24

METHODOLOGY	24

3.1  System Overview	24

3.1.1  Design Goals and Requirements	24

3.1.2  System Architecture Overview	25

3.1.3  Design Rationale: Why eBPF at the VFS Layer	27

3.1.4  Threat Model and Assumptions	28

3.2  Data Collection	30

3.2.1  Isolated Virtual Machine Environment	30

3.2.2  eBPF Sensor Design	31

3.2.3  Benign Workload Generator	32

3.2.4  Automated Orchestration Pipeline	32

3.2.5  CSV-to-Parquet Conversion and Labeling	33

3.3  Dataset	34

3.3.1  Ransomware Families	34

3.3.2  Dataset Composition	36

3.3.3  Labeling Methodology	37

3.3.4  Data Quality Audit	39

3.3.5  Class Imbalance	39

3.3.6  Noise Filtering	41

3.4  Tokenization Scheme	41

3.4.1  10-bit Token Layout	41

3.4.2  Effective Signal Analysis	44

3.5  Model Architecture	44

3.5.1  Stateful \textmu{}LSTM	44

3.5.2  Truncated Backpropagation Through Time	45

3.5.3  Hyperparameters	45

3.6  Training Methodology	46

3.6.1  Leave-One-Family-Out Protocol	46

3.6.2  3-Way Data Split	48

3.6.3  Loss Function and Class Weighting	48

3.6.4  Early Stopping and Learning Rate Schedule	49

3.7  Evaluation Methodology	49

3.7.1  Primary Metrics	49

3.7.2  Per-Family Evaluation	50

3.8  Real-Time Deployment Prototype	50

3.8.1  System Design	51

3.8.2  3-Layer Detection Pipeline	51

CHAPTER IV	54

EXPERIMENTAL RESULTS	54

4.1  Chapter Overview	54

4.2  Evaluation Setup	54

4.3  Leave-One-Family-Out (LOFO) Results	57

4.4  Per-Family Analysis	60

4.4.1  Behavioral Groups	60

4.4.2  Early-Detection Traces	63

4.5  Error Analysis	64

4.5.1  Log-Spam Dilution (Akira 0ee1d284ed663073)	64

4.5.2  Sensor Blind Spots (Guunra and Akira pwritev variants)	66

4.5.3  Root Causes of Token-Representation Failures	68

4.6 Heuristic False-Positive Suppression	69

4.7  Production Model Configuration and Inference Microbenchmark	71

4.8  Live-VM Validation	75

4.8.1  Motivation and Test Environment	75

4.8.2  Tests and Observations	75

4.8.3  Mixed-Workload Resilience	77

4.8.4  Limitations	78

4.9  Aggregate-Feature Baseline Comparison and Threats to Validity	79

4.10  Chapter Summary	81

CHAPTER V	83

CONCLUSION	83

5.1  Summary of Findings and Answers to the Research Questions	83

5.1.1  \textmu{}LSTM Classifier Detection	83

5.1.2  False-Positive Reduction	84

5.1.3  MBD3 and Data Loss Exposure	84

5.1.4  Mixed-Workload Feasibility	84

5.2  Comparison with the CLEAR Framework	85

5.3  Limitations	85

5.3.1  Dataset and Environment Scope	85

5.3.2  Sensor Visibility and OpCode Coverage	86

5.3.3  Metric Interpretation	86

5.3.4  Limitations Identified in Live Deployment	86

5.4  Future Work	86

5.4.1  Expanded Family Coverage	86

5.4.2  Production Deployment and Sensor Extensions	87

5.5  Concluding Remarks	87

REFERENCES	88

Appendix A Malware Sample Catalog	93

A.1  Sample-Acquisition Summary	93

A.2  Binary Samples Included in the Evaluation Dataset	93

A.3  Excluded Families	94

A.4  Reproducibility Notes	94

Appendix B Benign Workload Specification	96

B.1  File-Pool Configuration	96

B.2  Implemented Task Families	96

B.3  Scheduling and Orchestrator Invocation	97

B.4  Capture, Labeling, and Tokenization	97

B.5  Scope and Reproducibility Limitations	98

Appendix C Sensor, Tokenization, and Data-Conversion Specification	99

C.1  Data-Collection Sensor Implementation	99

C.1.1  Path Reconstruction and Filtering	99

C.1.2  Buffering and Probe Attachment	100

C.2  Event Tokenization	100

C.2.1  Offset and Auxiliary-State Semantics	100

C.2.2  Differences from CLEAR	100

C.2.3  Post-Collection Noise Filter	101

C.3  CSV-to-Parquet Conversion and Output Artifacts	101

C.4  Reproducibility Boundaries	102

Appendix D Run Orchestration and Data-Quality Audit Specification	103

D.1  Quality-Control Stages	103

D.2  Implemented Data-Quality Audit	103

D.2.1  Recordings Excluded after Audit	104

D.3  Orchestration Sequence and Timing	104

D.3.1  Duration Caps	105

D.3.2  Retry Hierarchy	105

D.4  Reproducibility Boundaries	106

Appendix E Real-Time Deployment Implementation Specification	107

E.1  Process Model and Thread Roles	107

E.2  Inter-Thread Communication	107

E.3  Deployment Decision Sequence	108

E.3.1  Heuristic Conditions	109

E.4  Buffering, Back-Pressure, and PID Lifecycle	110

E.5  Alert Actions and Shutdown	111

E.6  Deployment Variants and Reproducibility Boundaries	111

Appendix F Supplementary Inference and Tokenization Ablations	113

F.1  Interpretation and Scope	113

F.2  Output-Frequency Ablation	113

F.3  Decision-Aggregation Ablation	114

F.4  Cross-Chunk Statefulness Ablation	115

F.5  Token-2 Auxiliary-Bit Ablation	116

F.6  SougoLock Offset-Resolution Ablation	117

LIST OF ABBREVIATIONS	118

LIST OF TABLES

LIST OF FIGURES
